\documentclass[a4paper,12pt]{article}

\usepackage[T2A]{fontenc}
\usepackage[utf8]{inputenc}
\usepackage[russian]{babel}

\usepackage{amsmath, amssymb, amsthm, mathtools}
\usepackage{helvet}
\renewcommand{\familydefault}{\sfdefault}
\usepackage{icomma}
\usepackage{geometry}
\usepackage{microtype}
\usepackage{tikz}
\usepackage{tkz-euclide}
\usepackage{pgfplots}
\pgfplotsset{compat=1.18}
\usepackage{tabularx, booktabs, longtable}
\usepackage{enumitem}
\usepackage{hyperref}
\usepackage{parskip}
\usepackage{fancyhdr}
\usepackage{setspace}
\usepackage{xcolor}

\definecolor{accent}{HTML}{008080}
\definecolor{darkaccent}{HTML}{004D4D}
\definecolor{textgray}{HTML}{333333}

\geometry{left=2.5cm, right=2.5cm, top=2.5cm, bottom=2.5cm}
\onehalfspacing

\begin{document}

% Титульный лист
\begin{titlepage}
\centering
\vspace*{2cm}
{\Large\bfseries\color{darkaccent} Чек-лист}\par
\vspace{0.5cm}
{{large Тема: Алгебраические преобразования. Степени и корни}}\par
\vspace{2cm}
{\large Лёвин Артём Александрович}\par
{\large эксклюзивно для Яра}\par
\vspace{2cm}
{\today}\par
\vfill
\begin{tikzpicture}
\draw[color=accent, line width=1.5pt] 
(0,0) .. controls (2,1.5) and (4,-0.5) .. (6,0.5) .. controls (8,1.5) and (10,-0.5) .. (12,0);
\end{tikzpicture}
\vspace{1cm}
\end{titlepage}

\section*{Чек-лист: Алгебраические преобразования. Степени и корни}

\subsection*{Основные понятия}
\begin{itemize}[leftmargin=*]
\item \textbf{Свойства степеней}: $a^n \cdot a^k = a^{n+k}$, $\frac{a^n}{a^k} = a^{n-k}$, $(a^n)^k = a^{nk}$.
\item \textbf{Формулы сокращённого умножения}: 
$a^2 - b^2 = (a-b)(a+b)$, $(a\pm b)^2 = a^2 \pm 2ab + b^2$.
\item \textbf{Квадратный корень}: $\sqrt{a^2} = |a|$, $\sqrt{ab} = \sqrt{a}\cdot\sqrt{b}$ при $a,b \geq 0$.
\item \textbf{Сопряжённое выражение}: используется для избавления от иррациональности в знаменателе.
\end{itemize}

\subsection*{Методы упрощения выражений}
\begin{enumerate}[leftmargin=*]
\item Выносить общие множители за скобки.
\item Применять формулы сокращённого умножения.
\item Сокращать дроби после разложения на множители.
\item Избавляться от иррациональности в знаменателе домножением на сопряжённое.
\item Приводить подобные слагаемые.
\end{enumerate}

\subsection*{Типичные ошибки}
\begin{itemize}[leftmargin=*]
\item[$\times$] Путают сложение и умножение степеней: $a^n + a^k \neq a^{n+k}$.
\item[$\times$] Забывают модуль при извлечении корня: $\sqrt{x^2} = |x|$.
\item[$\times$] Не берут подмножаемые в скобки при подстановке отрицательных значений.
\item[$\times$] Сокращают слагаемые вместо множителей.
\item[$\times$] Неправильно работают с трёхэтажными дробями.
\end{itemize}

\newpage
\section*{Тренировочный блок}

\textbf{Уровень 1 (базовый)}

1. Упростите выражение: $\frac{a^7 \cdot a^3}{a^8}$.

2. Вычислите: $\frac{15(a^2+b^2)}{a^2-b^2}$ при $a=2\sqrt{3}$, $b=\sqrt{2}$.

3. Упростите: $\frac{x^4-y^4}{x-y} : (x+y)$.

\medskip
\textbf{Уровень 2 (средний)}

4. Упростите: $\frac{9b^2}{a-4} \cdot \frac{a^2-16}{3b}$.

5. Вычислите: $\left(\frac{3}{4} : \frac{5}{2}\right) \cdot 10$.

6. Упростите: $\frac{\sqrt{18}}{\sqrt{2}} + \sqrt{9}$.

\medskip
\textbf{Уровень 3 (выше среднего)}

7. Упростите: $\frac{16a^2-25b^2}{4a-5b}$.

8. Вычислите: $\frac{(2-\sqrt{3})^2 + 4\sqrt{3}}{\sqrt{4+2\sqrt{3}}}$.

9. Упростите: $\frac{p^4-q^4}{p^2-q^2} : \frac{p^2+q^2}{p+q}$.

\medskip
\textbf{Уровень 4 (сложный)}

10. Найдите значение: $\sqrt{11-6\sqrt{2}} + \sqrt{2}$.


\end{document}
